\section{Operating Boundaries and an Agent Tool Loop}
\label{app:operating-study}

This study tests the systems principle with a different, deliberately small
response family. Its purpose is to separate three questions: whether the
available instruments can justify the goal, whether the admitted model
supports the observations, and whether a more elaborate allocator is useful.
We also execute an actual language-model agent through the evidence interface.
All observations are generated locally; no physical wireless deployment is
claimed. The protocol, implementation, goals, sample sizes, and analysis were
hashed before evaluation. This is a local freeze, not an externally registered
protocol. We retain every development and evaluation attempt.

\emph{Prior-information correction.} The frozen results below use zero
lower bounds on full-stage means. Section~\ref{app:prior-sensitivity}
replays the one-counter comparison with known nominal radio service minima
supplied equally to both verifiers. Both then complete every trial at the
same cost; the reported joint-versus-componentwise separation therefore
does not persist with those physically justified priors. We retain the
original outcomes to make that sensitivity explicit.

\subsection{Population, Measurements, and Independent Evaluation}

There are two radio choices and two server choices. We allow the server
response to depend on the radio choice:
$$
\theta=(r_0,r_1,w_{00},w_{01},w_{10},w_{11}),
\qquad D_{ij}=r_i+w_{ij}.
$$
These are full stage means, including nominal service and transport terms.
Four probes measure aggregate response; two additional probes measure the
radio stages. Aggregate directions have rank four, aggregate directions with
one radio counter have rank five, and all six directions have rank six.
The radio response remains shared across server choices by construction.
Known configuration charges are $(1.6,2.4,2.6,3.4)$, with tolerance $0.05$.
The joint goal requires total, radio, and server means no larger than
$24$, $3.8$, and $8.5$ ms, respectively.

We reuse the frozen finite-buffer FCFS workload populations and independently
resample 64 virtual tasks per report. Queue populations record pre-admission
workload, including arrivals whose subsequent work is rejected. Full virtual
diagnostic response adds nominal stage terms. Thus the target is the mean of
this explicitly defined diagnostic population, not completion latency over
all offered jobs. Independent resampling justifies the report concentration
model; it does not establish independence for consecutive physical traffic.
The evaluator uses exact finite-population expectations, which never enter
the controller's public state. No separately estimated physical reference
mean is available or needed for this conditional simulator experiment.

Every epoch in studies A--C starts with 256 reports per aggregate probe and
no radio reports. We report their acquisition as sunk history. Online batches
contain 32 reports, and each original instrument count is capped at 4096.
For report variance proxy $\sigma_e^2$ and online error allocation
$\delta_{\rm on}$, the simultaneous radius is
$$
h_e(n)=\sqrt{\frac{2\sigma_e^2}{n}
                    \log\frac{2\cdot6\cdot4096}{\delta_{\rm on}}}.
$$
Known component workload caps supply the proxies, including the reduction
from averaging independent cohort components. We intersect every admitted
prefix interval with previous intervals, public mean bounds, and any
calibrated contrasts. A goal change does not reset this confidence ledger.
Empty intersections return model conflict and never authorize an action.

LP dual witnesses bound candidate constraints and exclude cheaper rivals.
We bound numerical stationarity residuals over the public variable box.
An optimizer-independent checker verifies the saved witness inequalities
with exact rational arithmetic applied to their binary floating-point inputs.
This validates the certificate algebra. The statistical coverage assumptions
and floating-point construction of confidence endpoints remain separate
obligations; this is not a machine-checked proof of the entire stochastic
implementation.

\subsection{Study A: Joint Certification and the Instrument Boundary}

We compare joint and sound componentwise verification on identical prefixes
of one fixed acquisition stream. It selects the available probe with the
smallest absolute report count; ties follow fast radio, slow radio, and then
the four aggregate probes. Instrument sets are aggregate-only, aggregate plus
the fast-radio counter, and all six probes. The 150- and 2000-unit duration
limits truncate this same stream at completed-batch boundaries. Each rule is
charged its own first certified prefix, or its complete unresolved prefix.
The common trace can continue solely to evaluate the other rule. This design
isolates the verifier, without claiming an optimal acquisition schedule.

The two states add either zero or 2 ms to every radio sample and subtract
the same amount from every full server sample. There is no clipping, and all
full stage samples remain nonnegative. Aggregate samples are bitwise identical
on matched tapes. Public bounds and initial history are also identical.
This shifted pair is a constructed bounded-stage family: subtracting from
full server response changes its original nominal service/transport baseline.
It is not presented as two realizations of an unchanged FCFS mechanism.
The unshifted state has acceptable answer $3$; the shifted state has no
feasible configuration. A uniformly $\delta$-correct aggregate-only procedure
therefore has correct-completion probability at most $\delta$ at either
state, by the indistinguishability argument. This is distinct from its
probability of any nonabstention, which can be as large as $2\delta$.

The joint advantage has a simple operational explanation. If an aggregate
lower bound exceeds $3.8+8.5=12.3$ ms, the two stage budgets cannot both hold.
The verifier can exclude that configuration without determining which stage
fails. It still needs enough stage evidence to establish the chosen
configuration's own feasibility. Componentwise exclusion instead requires a
particular stage to be uniformly invalid over the same confidence set.

We use 32 independent acquisition seeds per condition; the two states and
verifiers are paired, not independent additional replications.
Table~\ref{tab:v41-capability} reports the complete 2000-unit comparison.
The completion counts are identical at 150 units. Joint verification uses
32 fast-radio reports; componentwise verification with all instruments uses
32 reports from each radio counter. With one counter, componentwise
verification remains unresolved after spending 6745.92 online cost units at
the longer duration limit. With aggregate-only access, both rules spend
2449.92 and remain unresolved. These costs include unsuccessful attempts.
The common initial history costs 1920 units in every condition.

\begin{table}[t]
\centering
\caption{Study A under the original zero-lower-bound prior: results in
each matched state at the 2000-unit duration limit. Entries are completion count and mean new cost.
The same completion counts hold at 150 units.}
\label{tab:v41-capability}
\small
\begin{tabular}{@{}lrr@{}}
\toprule
Available instruments & Joint & Componentwise\\
\midrule
Aggregate only & $0/32$; 2449.92 & $0/32$; 2449.92\\
Aggregate + fast radio & $32/32$; 384 & $0/32$; 6745.92\\
All six & $32/32$; 384 & $32/32$; 768\\
\bottomrule
\end{tabular}
\end{table}

At 150 units, unresolved aggregate-only runs cost 169.92 and unresolved
single-counter componentwise runs cost 768. The 95\% Wilson intervals are
$[89.3,100]\%$ for $32/32$ completion and $[0,10.7]\%$ for $0/32$.
These intervals describe repeated report sampling from the fixed populations.
They do not measure variation across physical networks. No incorrect answer
is observed. The two tested duration limits yield no additional completion
transition; this study establishes the stated instrumentation and verifier
contrast, rather than locating a resource-optimal completion boundary.

\subsection{Study B: Interactions, Calibration, and Model Conflict}

This study keeps the instrumentation boundary separate from state-specific
stage calibration. We increase server load by the coupling coefficient times
the radio index, capped at $0.96$, for coefficients $0$ and $0.24$. The
two-radio scaling is explicitly different from the earlier three-radio
diagnostic's index divided by two. The same archived queue generator creates
the coupled populations. We compare a flexible six-mean family, a family
with independently calibrated signed contrasts, and the nominal restriction
$w_{1j}-w_{0j}=0$. The last is a potentially invalid negative control, not a
guaranteed comparator under interactions.

Calibration acquires 512 independent cohort reports from each of the four
server populations. Only the two contrast intervals leave the calibrator.
For each difference of independent server-report averages, its radius uses
the sum of the two bounded variance proxies and a union bound over the two
contrasts. Calibration receives error allocation $0.02$; online inference
receives $0.03$. Flexible and nominal runs use $0.05$ online. Every calibration
costs 24,576 acquisition units and 4096 modeled duration units. We retain all
32 calibrations, including their observed coverage status.

Operational server means then receive a common $0.4$-ms offset. This makes
calibration's absolute means non-current while preserving the contrasts.
Contrast transport is a declared generator property; the experiment does
not infer a validity horizon or certify unannounced drift. All three models
receive the same online report stream and a 2000-unit duration limit.

Without coupling, all three models complete all 16 decisions correctly.
Mean online cost is 2208 for the flexible and nominal models and 2448 for
the calibrated model. The latter's smaller online error allocation widens
its online bands, and its contrast information does not recover that cost
in this fixture. With coupling $0.24$, the flexible and calibrated models
correctly certify infeasibility from history in all 16 cases. The nominal
model reports conflict in every case; it makes no incorrect approval.
All 32 independently calibrated contrast intervals cover the population
contrasts. A $16/16$ completion fraction has Wilson interval
$[80.6,100]\%$; zero observed calibration failures is not a physical transfer
guarantee. Including its 24,576-unit calibration ledger, the calibrated
method provides no demonstrated cost advantage here. The flexible response
family already handles this declared interaction without calibration.

\subsection{Study C: Allocation with the Same Joint Verifier}

We separately compare coverage, the archived residual minimum-cost design,
and the archived cost-aware alternative-information maximin design. Each
coupling receives 16 new paired seeds, with all six instruments, the flexible
model, the joint goal, and a 2000-unit duration limit. These allocation
objectives use Gaussian quadratic information as a heuristic; the common
bounded-report verifier alone controls authorization. They are adaptations
of the specified archived routines, not reproductions of a published
networking algorithm with transferred guarantees.

The adapter uses regularization only to choose a planning center, never to
construct the certificate. If its initial center is infeasible, it searches
the public confidence set for a feasible configuration. A linear program
maximizes common constraint slack; a second finds the nearest center while
retaining one quarter of that slack. Both objectives use this same rule.
Plans are cached for at most 256 added reports. Maximin uses 12 iterations;
residual uses the stated confidence-based quadratic target. Every cap,
affordability failure, solver failure, and coverage fallback is recorded.
Development initially exposed a center-selection fallback; its original
outputs are preserved separately from the corrected, frozen implementation.

All methods return all 16 decisions correctly at each coupling.
Without coupling, coverage costs 384, maximin costs 454.08, and residual
costs 524.16. Coverage therefore uses 70.08 fewer cost units than maximin
and 140.16 fewer than residual. Each method takes the same number of
batches across its 16 seeds, so the paired bootstrap cost intervals collapse
to those respective differences. This reflects the fixed fixture and batch
resolution, not certainty about performance elsewhere. Both information
routines execute on all 16 nontrivial episodes without a coverage fallback.
Their mean full-controller times are 0.076 s for maximin and 0.143 s for
residual, compared with 0.038 s for coverage. These descriptive times include
simulation and verification and were not obtained in a randomized timing
benchmark. With coupling $0.24$, history already certifies infeasibility;
all methods cost zero and invoke no acquisition planner. This adverse result
supports simple allocation as a reasonable implementation choice. It remains
separate from the earlier frozen 2.36\% maximin and 6.74\% G-ACOL-adapter
comparisons.

\subsection{Study D: Actual Agent Integration and Announced Epoch Changes}

A fresh language-model agent receives only a public goal and evidence
projection and uses the documented tools to propose configurations, acquire
reports, verify answers, and authorize or return unresolved. It receives no
fixture means, random seeds, or outcome labels. This is protocol isolation,
not an operating-system access boundary. We retain the prompt, controller
metadata, and every action and public response. The execution interface does
not expose an exact model-version identifier, so this run cannot support a
version-specific replication or model comparison.

The four goals are loose total delay, the joint stage goal, loose total
delay again, and the same joint goal after an announced state change. The
first three share an epoch and evidence ledger. The state change quarantines
the old reports and starts fresh inference. Each epoch receives risk $0.025$
and separate cost and modeled-duration caps of 2000 units. Only the first
epoch starts with the sunk aggregate history. We also execute a deterministic
public-interface coverage controller on the same potential tapes. Invalid
proposals, unavailable probes, stale tokens, idempotent retries, and history
quarantine are exercised in separately labeled scripted guard tests.

The actual agent closes all four goals correctly: configuration $0$,
configuration $3$, configuration $0$, and certified infeasibility after the
announced change. It makes 23 accepted actions, no rejected requests, and
no unresolved terminal decisions. It acquires 64 fast-radio reports for the
joint goal. After the epoch change, it acquires 128 aggregate reports for
configuration 1 and 64 reports from each of aggregates 0, 2, 3 and the
fast-radio counter. Total new acquisition is 448 reports, costing 2133.76
units and 729.44 modeled duration units. The two loose goals require no new
reports, including the return to the loose goal after stage acquisition.

The deterministic control also closes all four goals correctly. It uses
45 accepted actions, costs 2124.80, and spends 735.20 modeled duration units.
Both controllers preserve evidence across goal changes and quarantine it at
the announced epoch transition. Their observed costs are similar; this one
functional sequence establishes neither an LLM performance advantage nor a
population completion rate. Sixteen separate scripted guard checks pass.
Independent loop replay passes 2944 cohort reports, 68 actions, and eight
correct authorized certificates across the two controllers. It reconstructs
evidence, token bindings, and epoch ledgers without importing the API or
verifier. The tool-operation time and transcript wall times are retained, with
orchestration delays identified; they are not isolated inference timings or
physical decision deadlines.

\subsection{Accounting, Audit, and What Remains Unresolved}

The release contains all attempts, raw acquired reports, terminal witnesses,
calibration ledgers, independent replay programs, and paired analyses.
Completion intervals are 95\% Wilson intervals. Cost differences use 10,000
paired seed-cluster bootstrap replicates. These intervals are descriptive;
we make no multiplicity-adjusted superiority claim. Common-stream wall time
includes evaluating both rules and is not isolated verifier latency.
Allocation-controller wall times are also recorded separately from abstract
probe durations. Logical simulator actuation has no modeled switching delay.
We therefore do not combine these quantities into a physical deadline claim.

The independent study auditor passes 320 recorded traces, 960 attempt
rows, 6598 raw batches containing 497,856 cohort reports, and 576 saved
algebraic certificates. It checks all 384 ordinary unresolved cost ledgers
and identifies the 16 nominal-model conflicts separately. There are 560
correct approvals and no observed wrong approvals across these correlated
factorial rows; this total is not an independent sample for a pooled error
rate interval. The auditor reconstructs tapes, calibration formulas,
confidence polytopes, resource accounting, and outcomes without importing
any experiment controller, verifier, or allocation adapter. It does not
establish that the selected allocation or stopping prefix is optimal.

The study supports a goal-dependent distinction between more reports and
another instrument, and it demonstrates a functioning verifier-controlled
agent interface. It does not establish a new general statistical result.
The two duration limits do not locate the optimal finite-resource frontier;
an observed unresolved prefix does not establish a budget impossibility.
Physical response validation, measured sensing and actuation times,
unannounced state-change handling, calibrated evidence lifetimes, larger
configuration sets, and tail-sensitive service specifications remain open.
The complete new study is under \texttt{research\_v41/study/}.
